% ch6-inner-products.tex — Chapter 6: Inner Product Spaces and Orthogonality
% Bilinear and sesquilinear forms, quadratic forms, inner products,
% norms, Cauchy–Schwarz, orthogonality, Gram–Schmidt, orthonormal bases,
% orthogonal/unitary matrices, adjoint operators, applications.

\chapter{Inner Product Spaces and Orthogonality}\label{ch:inner-products}

% ============================================================
% Motivating introduction
% ============================================================

Up to this point, our study of linear algebra has been purely
\emph{algebraic}: we have manipulated vector spaces, linear maps, and
matrices without ever measuring a length, an angle, or a distance.  Yet
geometry---in the intuitive sense of Euclidean space---depends critically
on such measurements.  How long is a vector?  What angle do two vectors
enclose?  Which vector in a subspace is closest to a given point?

The tool that answers all these questions at once is the \emph{inner
product}\index{inner product}: a function that takes two vectors and
returns a scalar, generalising the familiar dot product on~$\R^n$.  Once
an inner product is fixed, the underlying vector space acquires a rich
geometric structure---norms, distances, angles, orthogonality,
projections---and many powerful theorems become available.

\begin{center}
\begin{tikzpicture}[>=Stealth, scale=1.6]
  % --- Angle between two vectors ---
  \draw[gray!30, very thin, step=0.5] (-0.3,-0.3) grid (3.8,2.8);
  \draw[->] (-0.3,0) -- (3.8,0) node[right]{$x_1$};
  \draw[->] (0,-0.3) -- (0,2.8) node[above]{$x_2$};

  \draw[->, thick, blue!70!black] (0,0) -- (3,1)
    node[right] {$\vec{u}$};
  \draw[->, thick, red!70!black] (0,0) -- (1,2.5)
    node[above] {$\vec{v}$};

  % Angle arc
  \draw[thick, purple!70!black] (0.9,0.3) arc[start angle=18.4,
    end angle=68.2, radius=0.95];
  \node[purple!70!black] at (0.75,0.95) {$\theta$};

  % Annotation
  \node[below, text width=4cm, align=center] at (1.8,-0.6)
    {\small $\cos\theta = \dfrac{\inner{\vec{u},\vec{v}}}%
      {\norm{\vec{u}}\;\norm{\vec{v}}}$};
\end{tikzpicture}
\captionof{figure}{The inner product encodes the angle between two
  vectors.  All of Euclidean geometry follows.}
\label{fig:inner-product-angle}
\end{center}

This chapter develops the theory systematically, starting from the
general notion of bilinear forms, specialising to inner products, and
building up the machinery of orthogonality, projections, and
orthonormal bases.  We then study the linear maps that preserve inner
products (orthogonal and unitary matrices) and those that are
``self-dual'' with respect to the inner product (adjoint operators).

Throughout, $\K$ denotes $\R$ or~$\C$, and $E$ denotes a
finite-dimensional $\K$-vector space unless stated otherwise.  We use
the convention that inner products on~$\C$ are linear in the
\emph{first} argument and conjugate-linear in the second (the
\emph{physics convention}).


% ============================================================
\section{Bilinear forms}\label{sec:bilinear-forms}
% ============================================================

\begin{definition}{Bilinear form}{def:bilinear-form}
\index{bilinear form}
Let $E$ be a vector space over~$\K$.  A \emph{bilinear form} on~$E$ is
a map $\varphi\colon E\times E \to \K$ that is linear in each argument:
\begin{enumerate}[label=(\roman*)]
  \item For all $u,v,w\in E$ and $\lambda\in\K$:
    \[
      \varphi(u+\lambda v,\, w) = \varphi(u,w) + \lambda\,\varphi(v,w).
    \]
  \item For all $u,v,w\in E$ and $\lambda\in\K$:
    \[
      \varphi(u,\, v+\lambda w) = \varphi(u,v) + \lambda\,\varphi(u,w).
    \]
\end{enumerate}
\end{definition}

\begin{definition}{Symmetric and antisymmetric bilinear forms}%
{def:symmetric-bilinear}
\index{bilinear form!symmetric}\index{bilinear form!antisymmetric}
A bilinear form $\varphi$ on~$E$ is called:
\begin{itemize}
  \item \emph{symmetric} if $\varphi(u,v) = \varphi(v,u)$ for all
    $u,v\in E$;
  \item \emph{antisymmetric} (or \emph{skew-symmetric}) if
    $\varphi(u,v) = -\varphi(v,u)$ for all $u,v\in E$.
\end{itemize}
\end{definition}

\begin{definition}{Matrix of a bilinear form}{def:matrix-bilinear}
\index{bilinear form!matrix of}
Let $\cB = (\vece{1},\ldots,\vece{n})$ be a basis of~$E$ and let
$\varphi$ be a bilinear form on~$E$.  The \emph{matrix of~$\varphi$
in the basis~$\cB$} is the matrix $G\in\Mat_n(\K)$ defined by
\[
  G_{ij} \deff \varphi(\vece{i},\vece{j}),
  \qquad 1\leq i,j\leq n.
\]
If $u = \sum_i x_i\,\vece{i}$ and $v = \sum_j y_j\,\vece{j}$, then
\[
  \varphi(u,v) = \transpose{X}\, G\, Y,
\]
where $X = (x_1,\ldots,x_n)^{\mathsf{T}}$ and
$Y = (y_1,\ldots,y_n)^{\mathsf{T}}$ are the coordinate vectors.
\end{definition}

\begin{proposition}{Change of basis for bilinear forms}%
{prop:bilinear-change-basis}
\index{bilinear form!change of basis}
If $P$ is the change-of-basis matrix from $\cB$ to $\cB'$, and $G$ is
the matrix of~$\varphi$ in~$\cB$, then the matrix of~$\varphi$
in~$\cB'$ is
\[
  G' = \transpose{P}\, G\, P.
\]
In particular, $\varphi$ is symmetric if and only if $G$ is symmetric
($G = \transpose{G}$).
\end{proposition}

\begin{proof}
Let $u,v\in E$ with coordinates $X'$ in~$\cB'$ and $X = PX'$ in~$\cB$
(similarly $Y = PY'$).  Then
\[
  \varphi(u,v) = \transpose{X}\, G\, Y
    = \transpose{(PX')}\, G\, (PY')
    = \transpose{X'}\,(\transpose{P}\, G\, P)\, Y'.
\]
Since this holds for all $X', Y'$, we have $G' = \transpose{P}\, G\, P$.
\end{proof}

\begin{definition}{Non-degenerate bilinear form}{def:non-degenerate}
\index{bilinear form!non-degenerate}\index{non-degenerate}
A bilinear form $\varphi$ on~$E$ is \emph{non-degenerate} if
\[
  \bigl(\,\forall v\in E,\;\varphi(u,v) = 0\bigr)
  \;\Longrightarrow\; u = \vzero.
\]
Equivalently, the matrix~$G$ of~$\varphi$ in any basis is invertible.
\end{definition}


% ============================================================
\section{Quadratic forms}\label{sec:quadratic-forms}
% ============================================================

\begin{definition}{Quadratic form}{def:quadratic-form}
\index{quadratic form}
Let $E$ be a $\K$-vector space.  A \emph{quadratic form} on~$E$ is a
map $q\colon E\to\K$ such that
\begin{enumerate}[label=(\roman*)]
  \item $q(\lambda v) = \lambda^2\, q(v)$ for all $\lambda\in\K$,
    $v\in E$;
  \item the map $\varphi\colon E\times E\to\K$ defined by
    \[
      \varphi(u,v) \deff \frac{1}{2}\bigl[q(u+v) - q(u) - q(v)\bigr]
    \]
    is bilinear (this is called the \emph{polar form}\index{polar form}
    of~$q$).
\end{enumerate}
Conversely, every symmetric bilinear form $\varphi$ defines a quadratic
form $q(v) = \varphi(v,v)$.
\end{definition}

\begin{proposition}{Polarization identity}{prop:polarization}
\index{polarization identity}
If $\Char(\K)\neq 2$ and $\varphi$ is a symmetric bilinear form with
associated quadratic form $q(v) = \varphi(v,v)$, then
\[
  \varphi(u,v) = \frac{1}{2}\bigl[q(u+v) - q(u) - q(v)\bigr]
    = \frac{1}{4}\bigl[q(u+v) - q(u-v)\bigr].
\]
Thus, $q$ and $\varphi$ determine each other uniquely.
\end{proposition}

\begin{proof}
Expanding $q(u+v) = \varphi(u+v,\,u+v) = q(u) + 2\varphi(u,v) + q(v)$
gives the first identity.  For the second, note that
$q(u-v) = q(u) - 2\varphi(u,v) + q(v)$, so subtracting yields
$q(u+v) - q(u-v) = 4\varphi(u,v)$.
\end{proof}

\begin{theorem}{Sylvester's law of inertia}{thm:sylvester}
\index{Sylvester's law of inertia}\index{signature}
Let $q$ be a quadratic form on a real vector space~$E$ of
dimension~$n$.  There exists a basis of~$E$ in which the matrix of the
associated bilinear form is diagonal with entries $+1$, $-1$, and~$0$:
\[
  G = \diag(\underbrace{1,\ldots,1}_{p},\;
    \underbrace{-1,\ldots,-1}_{r-p},\;
    \underbrace{0,\ldots,0}_{n-r}).
\]
The integers $p$ (the number of $+1$'s), $r-p$ (the number of $-1$'s),
and $n-r$ (the number of~$0$'s) depend only on~$q$, not on the choice
of basis.  The pair $(p,\,r-p)$ is called the \emph{signature} of~$q$,
and $r = \rank(q)$.
\end{theorem}

\begin{remark}{Positive definiteness from the signature}%
{rem:signature-positive-definite}
A real quadratic form~$q$ is \emph{positive definite}\index{positive
definite} (i.e.\ $q(v)>0$ for all $v\neq\vzero$) if and only if its
signature is $(n,0)$, i.e.\ all signs are~$+1$.
\end{remark}


% ============================================================
\section{Inner products}\label{sec:inner-products}
% ============================================================

\begin{definition}{Inner product (real case)}{def:inner-product-real}
\index{inner product!real}
Let $E$ be a real vector space.  An \emph{inner product} on~$E$ is a
bilinear form $\inner{\cdot,\cdot}\colon E\times E\to\R$ that is
\begin{enumerate}[label=(\roman*)]
  \item \emph{symmetric}: $\inner{u,v} = \inner{v,u}$ for all
    $u,v\in E$;
  \item \emph{positive definite}: $\inner{v,v} > 0$ for all
    $v\neq\vzero$.
\end{enumerate}
A real vector space equipped with an inner product is called a
\emph{Euclidean space}\index{Euclidean space}.
\end{definition}

\begin{definition}{Sesquilinear form}{def:sesquilinear}
\index{sesquilinear form}
Let $E$ be a complex vector space.  A \emph{sesquilinear form} on~$E$
is a map $\varphi\colon E\times E\to\C$ that is linear in the first
argument and conjugate-linear in the second:
\begin{enumerate}[label=(\roman*)]
  \item $\varphi(\alpha u + \beta v,\, w) = \alpha\,\varphi(u,w)
    + \beta\,\varphi(v,w)$;
  \item $\varphi(u,\, \alpha v + \beta w) = \bar{\alpha}\,\varphi(u,v)
    + \bar{\beta}\,\varphi(u,w)$.
\end{enumerate}
A sesquilinear form is \emph{Hermitian}\index{Hermitian form} if
$\varphi(u,v) = \overline{\varphi(v,u)}$ for all $u,v\in E$.
\end{definition}

\begin{definition}{Inner product (complex case)}{def:inner-product-complex}
\index{inner product!complex}\index{Hermitian inner product}
Let $E$ be a complex vector space.  A \emph{Hermitian inner product}
(or simply \emph{inner product}) on~$E$ is a sesquilinear form
$\inner{\cdot,\cdot}\colon E\times E\to\C$ that is
\begin{enumerate}[label=(\roman*)]
  \item \emph{Hermitian symmetric}: $\inner{u,v} =
    \overline{\inner{v,u}}$ for all $u,v\in E$;
  \item \emph{positive definite}: $\inner{v,v} > 0$ for all
    $v\neq\vzero$ (note that $\inner{v,v}\in\R$ by~(i)).
\end{enumerate}
A complex vector space equipped with a Hermitian inner product is called
a \emph{unitary space}\index{unitary space} (or \emph{pre-Hilbert
space}).
\end{definition}

\begin{remark}{Unified notation}{rem:unified-inner-product}
Over~$\R$, a Hermitian inner product reduces to a symmetric bilinear
form (since conjugation is trivial).  We therefore write
$\inner{u,v}$ in both cases and speak simply of an ``inner product
space''\index{inner product space}.
\end{remark}


% ============================================================
\section{Examples}\label{sec:inner-product-examples}
% ============================================================

\begin{example}{Standard dot product on $\R^n$}{ex:dot-product}
\index{dot product}\index{standard inner product!on $\R^n$}
The \emph{standard inner product} on $\R^n$ is
\[
  \inner{x,y} \deff \sum_{i=1}^n x_i\, y_i = \transpose{x}\, y,
  \qquad x,y\in\R^n.
\]
Its matrix in the standard basis is $G = I_n$.
\end{example}

\begin{example}{Hermitian product on $\C^n$}{ex:hermitian-Cn}
\index{standard inner product!on $\C^n$}
The \emph{standard Hermitian inner product} on $\C^n$ is
\[
  \inner{z,w} \deff \sum_{i=1}^n z_i\,\bar{w}_i = \hermit{w}\, z,
  \qquad z,w\in\C^n.
\]
This is linear in the first argument and conjugate-linear in the second.
\end{example}

\begin{example}{$L^2$ inner product on $\mathcal{C}([a,b])$}%
{ex:L2-inner-product}
\index{L2 inner product@$L^2$ inner product}
On the space $\mathcal{C}([a,b],\R)$ of continuous real-valued functions
on~$[a,b]$, the map
\[
  \inner{f,g} \deff \int_a^b f(t)\, g(t)\,\mathrm{d}t
\]
is an inner product.  Symmetry and bilinearity are immediate from the
properties of the integral; positive definiteness follows because
$\inner{f,f} = \int_a^b f(t)^2\,\mathrm{d}t \geq 0$, with equality if
and only if $f = 0$ (by continuity).

The complex analogue on $\mathcal{C}([a,b],\C)$ is
$\inner{f,g} = \int_a^b f(t)\,\overline{g(t)}\,\mathrm{d}t$.
\end{example}

\begin{example}{Weighted inner product}{ex:weighted-inner-product}
\index{inner product!weighted}
Let $w_1,\ldots,w_n > 0$ be positive reals.  Then
\[
  \inner{x,y}_w \deff \sum_{i=1}^n w_i\, x_i\, y_i
\]
defines an inner product on~$\R^n$.  Its matrix in the standard basis is
$G = \diag(w_1,\ldots,w_n)$.
\end{example}

\begin{example}{Frobenius inner product}{ex:frobenius}
\index{Frobenius inner product}
On the space $\Mat_n(\R)$, the map
\[
  \inner{A,B} \deff \tr(\transpose{A}\, B) = \sum_{i,j} a_{ij}\, b_{ij}
\]
defines an inner product (the \emph{Frobenius inner product}).
\end{example}


% ============================================================
\section{Norm and distance}\label{sec:norm-distance}
% ============================================================

\begin{definition}{Norm induced by an inner product}{def:norm}
\index{norm}\index{norm!induced by inner product}
Let $(E,\inner{\cdot,\cdot})$ be an inner product space.  The
\emph{norm} of $v\in E$ is
\[
  \norm{v} \deff \sqrt{\inner{v,v}}.
\]
\end{definition}

\begin{definition}{Distance}{def:distance}
\index{distance}
The \emph{distance} between $u$ and~$v$ in an inner product space is
\[
  d(u,v) \deff \norm{u - v}.
\]
\end{definition}

\begin{proposition}{Basic properties of the norm}{prop:norm-properties}
Let $(E,\inner{\cdot,\cdot})$ be an inner product space.  For all
$u,v\in E$ and $\lambda\in\K$:
\begin{enumerate}[label=(\roman*)]
  \item $\norm{v}\geq 0$, with equality if and only if $v = \vzero$;
  \item $\norm{\lambda v} = \abs{\lambda}\,\norm{v}$ \quad
    (absolute homogeneity);
  \item $\norm{u+v}^2 + \norm{u-v}^2 = 2\norm{u}^2 + 2\norm{v}^2$
    \quad (parallelogram law\index{parallelogram law}).
\end{enumerate}
\end{proposition}

\begin{proof}
Parts (i) and (ii) follow directly from the definition and the
properties of the inner product.  For~(iii), expand using bilinearity:
\[
  \norm{u+v}^2 + \norm{u-v}^2
    = \inner{u+v,u+v} + \inner{u-v,u-v}
    = 2\inner{u,u} + 2\inner{v,v}
    = 2\norm{u}^2 + 2\norm{v}^2. \qedhere
\]
\end{proof}


% ============================================================
\section{The Cauchy--Schwarz inequality}\label{sec:cauchy-schwarz}
% ============================================================

\begin{theorem}{Cauchy--Schwarz inequality}{thm:cauchy-schwarz}
\index{Cauchy--Schwarz inequality}
Let $(E,\inner{\cdot,\cdot})$ be an inner product space over~$\K$.
For all $u,v\in E$,
\[
  \abs{\inner{u,v}} \leq \norm{u}\,\norm{v},
\]
with equality if and only if $u$ and~$v$ are linearly dependent
(i.e.\ one is a scalar multiple of the other).
\end{theorem}

\begin{proof}
If $v = \vzero$, both sides are zero and the result is trivial.  Assume
$v\neq\vzero$.

\medskip\noindent\textbf{Real case ($\K = \R$).}\quad
For any $t\in\R$, the positive definiteness of the inner product gives
\[
  0 \leq \norm{u - tv}^2
    = \inner{u - tv,\, u - tv}
    = \norm{u}^2 - 2t\inner{u,v} + t^2\norm{v}^2.
\]
This is a quadratic polynomial in~$t$ that is non-negative for all~$t$.
A real quadratic $at^2 + bt + c\geq 0$ for all~$t$ requires the
discriminant to satisfy $b^2 - 4ac \leq 0$.  Here
$a = \norm{v}^2$, $b = -2\inner{u,v}$, $c = \norm{u}^2$, so
\[
  4\inner{u,v}^2 - 4\norm{u}^2\norm{v}^2 \leq 0,
\]
which gives $\inner{u,v}^2 \leq \norm{u}^2\norm{v}^2$, and taking
square roots yields $\abs{\inner{u,v}}\leq\norm{u}\,\norm{v}$.

Equality holds if and only if the discriminant is zero, i.e.\ the
minimum of the quadratic is zero, which occurs at
$t_0 = \inner{u,v}/\norm{v}^2$, giving $u = t_0 v$.

\medskip\noindent\textbf{Complex case ($\K = \C$).}\quad
For any $t\in\R$ and any $\alpha\in\C$ with $\abs{\alpha}=1$, consider
\[
  0 \leq \norm{u - t\alpha v}^2
    = \norm{u}^2 - 2t\,\Re(\alpha\,\inner{v,u})
      + t^2\norm{v}^2.
\]
Choose $\alpha$ so that $\alpha\,\inner{v,u} = \abs{\inner{u,v}}$
(this is possible by writing $\inner{u,v} = \abs{\inner{u,v}}\,e^{i\theta}$
and setting $\alpha = e^{-i\theta}$).  Then
$\Re(\alpha\,\inner{v,u}) = \abs{\inner{u,v}}$, and the above becomes
\[
  0 \leq \norm{u}^2 - 2t\abs{\inner{u,v}} + t^2\norm{v}^2.
\]
This is again a non-negative real quadratic in~$t$, and the discriminant
argument gives
\[
  4\abs{\inner{u,v}}^2 - 4\norm{u}^2\norm{v}^2 \leq 0,
\]
i.e.\ $\abs{\inner{u,v}}\leq\norm{u}\,\norm{v}$.  Equality analysis is
identical: it holds if and only if $u = t_0\alpha v$ for the optimal
$t_0$, i.e.\ $u$ and~$v$ are linearly dependent.
\end{proof}

\begin{corollary}{Triangle inequality}{cor:triangle-inequality}
\index{triangle inequality}
For all $u,v\in E$,
\[
  \norm{u+v}\leq\norm{u}+\norm{v}.
\]
In particular, $d(\cdot,\cdot)$ satisfies the triangle inequality and
defines a metric on~$E$.
\end{corollary}

\begin{proof}
We compute
\begin{align*}
  \norm{u+v}^2
    &= \inner{u+v,\,u+v}
    = \norm{u}^2 + 2\Re\inner{u,v} + \norm{v}^2 \\
    &\leq \norm{u}^2 + 2\abs{\inner{u,v}} + \norm{v}^2 \\
    &\leq \norm{u}^2 + 2\norm{u}\,\norm{v} + \norm{v}^2
    = (\norm{u} + \norm{v})^2,
\end{align*}
where the last inequality uses Cauchy--Schwarz.  Taking square roots
gives the result.
\end{proof}

\begin{definition}{Angle between vectors}{def:angle}
\index{angle between vectors}
In a real inner product space, for nonzero $u,v\in E$, the
\emph{angle}~$\theta\in[0,\pi]$ between $u$ and~$v$ is defined by
\[
  \cos\theta \deff \frac{\inner{u,v}}{\norm{u}\,\norm{v}}.
\]
This is well-defined by the Cauchy--Schwarz inequality, which ensures
$\abs{\cos\theta}\leq 1$.
\end{definition}


% ============================================================
\section{Orthogonality}\label{sec:orthogonality}
% ============================================================

\begin{definition}{Orthogonal vectors}{def:orthogonal-vectors}
\index{orthogonal vectors}
Two vectors $u,v$ in an inner product space are \emph{orthogonal},
written $u\perp v$, if $\inner{u,v} = 0$.
\end{definition}

\begin{proposition}{Pythagorean theorem}{prop:pythagorean}
\index{Pythagorean theorem}
If $u\perp v$, then $\norm{u+v}^2 = \norm{u}^2 + \norm{v}^2$.
\end{proposition}

\begin{proof}
$\norm{u+v}^2 = \norm{u}^2 + 2\Re\inner{u,v} + \norm{v}^2
= \norm{u}^2 + \norm{v}^2$, since $\inner{u,v} = 0$.
\end{proof}

\begin{definition}{Orthogonal complement}{def:orthogonal-complement}
\index{orthogonal complement}
Let $F$ be a subset (or subspace) of an inner product space~$E$.  The
\emph{orthogonal complement} of~$F$ is
\[
  F^{\perp} \deff \setcond{v\in E}{\inner{v,w} = 0
    \text{ for all } w\in F}.
\]
\end{definition}

\begin{theorem}{Properties of the orthogonal complement}%
{thm:orth-complement-properties}
\index{orthogonal complement!properties}
Let $(E,\inner{\cdot,\cdot})$ be a finite-dimensional inner product
space and let $F$ be a subspace of~$E$.  Then:
\begin{enumerate}[label=(\roman*)]
  \item $F^{\perp}$ is a subspace of~$E$.
  \item $E = F \oplus F^{\perp}$ (orthogonal direct sum).
  \item $\dim F^{\perp} = \dim E - \dim F$.
  \item $(F^{\perp})^{\perp} = F$.
  \item $F\cap F^{\perp} = \set{\vzero}$.
\end{enumerate}
\end{theorem}

\begin{proof}
\textbf{(i)} $F^{\perp}$ is clearly non-empty ($\vzero\in F^{\perp}$)
and closed under addition and scalar multiplication by linearity of the
inner product.

\textbf{(v)} If $v\in F\cap F^{\perp}$, then $\inner{v,v} = 0$, so
$v = \vzero$ by positive definiteness.

\textbf{(ii)--(iv)} We prove these after establishing orthogonal
projection (
ef{thm:thm:orthogonal-projection}).  Once we know that every
$v\in E$ can be written as $v = p + p^{\perp}$ with $p\in F$ and
$p^{\perp}\in F^{\perp}$, the direct sum decomposition follows
from~(v).  Part~(iii) then follows by counting dimensions, and~(iv)
is immediate from $E = F\oplus F^{\perp}$ and the observation that
$F\subset (F^{\perp})^{\perp}$ together with the dimension count
$\dim(F^{\perp})^{\perp} = n - \dim F^{\perp} = \dim F$.
\end{proof}

\begin{proposition}{Orthogonal complement of sums and intersections}%
{prop:orth-sum-intersection}
Let $F_1,F_2$ be subspaces of a finite-dimensional inner product
space~$E$.  Then:
\begin{enumerate}[label=(\roman*)]
  \item $(F_1 + F_2)^{\perp} = F_1^{\perp}\cap F_2^{\perp}$.
  \item $(F_1\cap F_2)^{\perp} = F_1^{\perp} + F_2^{\perp}$.
\end{enumerate}
\end{proposition}

\begin{proof}
\textbf{(i)} $v\in(F_1+F_2)^{\perp}$ iff $\inner{v,w}=0$ for all
$w\in F_1+F_2$, which (since $F_1,F_2\subset F_1+F_2$) is equivalent
to $v\in F_1^{\perp}$ and $v\in F_2^{\perp}$.

\textbf{(ii)} Take orthogonal complements in (i) and use
$(F^{\perp})^{\perp}=F$.
\end{proof}


% ============================================================
\section{Orthogonal projection}\label{sec:orthogonal-projection}
% ============================================================

\begin{theorem}{Orthogonal projection}{thm:orthogonal-projection}
\index{orthogonal projection}
Let $(E,\inner{\cdot,\cdot})$ be a finite-dimensional inner product
space and let $F$ be a subspace of~$E$.  For every $v\in E$, there
exists a unique vector $p\in F$ such that
\[
  v - p \in F^{\perp}.
\]
The vector~$p$ is called the \emph{orthogonal projection} of~$v$
onto~$F$, denoted $p = \proj_F(v)$.  The map
$\proj_F\colon E\to E$ is a linear map satisfying
$\proj_F^2 = \proj_F$, $\im(\proj_F) = F$, and
$\Ker(\proj_F) = F^{\perp}$.
\end{theorem}

\begin{proof}
\textbf{Existence.}\quad Let $(e_1,\ldots,e_k)$ be an orthonormal basis
of~$F$ (which exists by Gram--Schmidt, 
ef{thm:thm:gram-schmidt}).  Set
\[
  p \deff \sum_{i=1}^k \inner{v,e_i}\, e_i.
\]
Then $p\in F$, and for each basis vector $e_j$ we have
\[
  \inner{v - p,\, e_j}
    = \inner{v,e_j} - \sum_{i=1}^k \inner{v,e_i}\inner{e_i,e_j}
    = \inner{v,e_j} - \inner{v,e_j} = 0.
\]
By linearity, $\inner{v-p,\,w} = 0$ for all $w\in F$, so
$v-p\in F^{\perp}$.

\textbf{Uniqueness.}\quad If $p'$ is another such vector, then
$p - p'\in F$ and $p - p'\in F^{\perp}$, so
$p-p'\in F\cap F^{\perp} = \set{\vzero}$.

\textbf{Linearity} is clear from the formula.  The remaining properties
$\proj_F^2 = \proj_F$, $\im(\proj_F) = F$, and
$\Ker(\proj_F) = F^{\perp}$ follow directly from the construction and
uniqueness.
\end{proof}

\begin{theorem}{Best approximation}{thm:best-approximation}
\index{best approximation}\index{orthogonal projection!best approximation}
Let $F$ be a subspace of a finite-dimensional inner product space~$E$
and let $v\in E$.  The orthogonal projection $p = \proj_F(v)$ is the
unique element of~$F$ closest to~$v$:
\[
  \norm{v - p} < \norm{v - w}
  \qquad\text{for all } w\in F,\; w\neq p.
\]
\end{theorem}

\begin{proof}
Let $w\in F$.  Since $v - p\in F^{\perp}$ and $p - w\in F$, the
Pythagorean theorem gives
\[
  \norm{v - w}^2
    = \norm{(v - p) + (p - w)}^2
    = \norm{v - p}^2 + \norm{p - w}^2
    \geq \norm{v - p}^2,
\]
with equality if and only if $p = w$.
\end{proof}

\begin{center}
\begin{tikzpicture}[>=Stealth, scale=1.5]
  % Subspace F (a plane, drawn as a parallelogram)
  \fill[blue!8] (-0.5,-0.3) -- (3.5,-0.3) -- (4.2,1.2) -- (0.2,1.2)
    -- cycle;
  \node[blue!70!black] at (3.8,0.15) {$F$};

  % Vector v
  \coordinate (v) at (1.8,2.5);
  \coordinate (p) at (1.8,0.5);

  \draw[->, thick, red!70!black] (0,0) -- (v)
    node[above right] {$v$};
  \draw[->, thick, blue!70!black] (0,0) -- (p)
    node[below right] {$p = \proj_F(v)$};

  % Perpendicular line
  \draw[thick, dashed, purple!70!black] (v) -- (p)
    node[midway, right, xshift=3pt] {$v - p \in F^{\perp}$};

  % Right angle mark
  \draw[purple!70!black] (1.8,0.75) -- (2.05,0.75) -- (2.05,0.5);

  % Another point w in F
  \coordinate (w) at (3,0.8);
  \draw[->, thick, gray] (0,0) -- (w) node[right] {$w$};
  \draw[dotted, thick, gray] (v) -- (w)
    node[midway, right, xshift=2pt] {\small $\norm{v-w}\geq\norm{v-p}$};
\end{tikzpicture}
\captionof{figure}{Orthogonal projection onto a subspace~$F$.  The
  projection $p = \proj_F(v)$ is the closest point in~$F$ to~$v$.}
\label{fig:orthogonal-projection}
\end{center}

\begin{proposition}{Projection formula with an orthonormal basis}%
{prop:projection-formula}
\index{orthogonal projection!formula}
If $(e_1,\ldots,e_k)$ is an orthonormal basis of a subspace~$F$, then
\[
  \proj_F(v) = \sum_{i=1}^k \inner{v,e_i}\, e_i.
\]
In particular, if $(e_1,\ldots,e_n)$ is an orthonormal basis of the
full space~$E$, then $v = \sum_{i=1}^n \inner{v,e_i}\, e_i$.
\end{proposition}


% ============================================================
\section{Gram--Schmidt orthogonalization}\label{sec:gram-schmidt}
% ============================================================

\begin{theorem}{Gram--Schmidt process}{thm:gram-schmidt}
\index{Gram--Schmidt process}\index{orthogonalization}
Let $(v_1,\ldots,v_k)$ be a linearly independent family in an inner
product space~$E$.  There exists a unique orthonormal family
$(e_1,\ldots,e_k)$ such that for each $j=1,\ldots,k$:
\begin{enumerate}[label=(\roman*)]
  \item $\Span(e_1,\ldots,e_j) = \Span(v_1,\ldots,v_j)$;
  \item $\inner{v_j,e_j} > 0$ (i.e.\ the ``leading coefficient'' is
    positive).
\end{enumerate}
The family $(e_1,\ldots,e_k)$ is constructed by the
\textbf{Gram--Schmidt algorithm}:
\begin{align*}
  w_1 &\deff v_1, &
    e_1 &\deff \frac{w_1}{\norm{w_1}}, \\
  w_j &\deff v_j - \sum_{i=1}^{j-1} \inner{v_j,e_i}\, e_i, &
    e_j &\deff \frac{w_j}{\norm{w_j}},
    \qquad j = 2,\ldots,k.
\end{align*}
\end{theorem}

\begin{proof}
We proceed by induction on~$j$.

\textbf{Base case ($j=1$).}\quad Set $w_1 = v_1$.  Since $v_1\neq\vzero$
(linear independence), $\norm{w_1}>0$ and $e_1 = w_1/\norm{w_1}$ is
well-defined with $\norm{e_1} = 1$.  Clearly
$\Span(e_1) = \Span(v_1)$ and $\inner{v_1,e_1} = \norm{v_1} > 0$.

\textbf{Inductive step.}\quad Assume that $(e_1,\ldots,e_{j-1})$ is an
orthonormal family with $\Span(e_1,\ldots,e_{j-1}) = \Span(v_1,\ldots,
v_{j-1})$ and $\inner{v_i,e_i}>0$ for $i<j$.  Define
\[
  w_j = v_j - \sum_{i=1}^{j-1}\inner{v_j,e_i}\, e_i.
\]
This is exactly $v_j - \proj_{F_{j-1}}(v_j)$, where
$F_{j-1} = \Span(e_1,\ldots,e_{j-1})$.  Hence $w_j\perp e_i$ for all
$i<j$ (by construction of the orthogonal projection).

We must show $w_j\neq\vzero$.  If $w_j = \vzero$, then
$v_j = \sum_{i=1}^{j-1}\inner{v_j,e_i}\,e_i\in F_{j-1} =
\Span(v_1,\ldots,v_{j-1})$, contradicting the linear independence
of $(v_1,\ldots,v_k)$.

Set $e_j = w_j/\norm{w_j}$.  Then $\norm{e_j} = 1$, and $e_j\perp e_i$
for all $i<j$ since $w_j\perp e_i$.  Moreover,
\[
  \Span(e_1,\ldots,e_j) = \Span(e_1,\ldots,e_{j-1},\,w_j)
    = \Span(e_1,\ldots,e_{j-1},\,v_j)
    = \Span(v_1,\ldots,v_j),
\]
using the induction hypothesis and the definition of~$w_j$.  Finally,
$\inner{v_j,e_j} = \inner{v_j, w_j/\norm{w_j}} = \norm{w_j} > 0$
(since $\inner{v_j,w_j} = \inner{w_j + \sum_i c_i e_i,\, w_j} =
\norm{w_j}^2 > 0$).

\textbf{Uniqueness.}\quad Suppose $(\tilde{e}_1,\ldots,\tilde{e}_k)$
also satisfies conditions (i)--(ii).  By~(i),
$\Span(\tilde{e}_1) = \Span(v_1)$, so $\tilde{e}_1 = \pm e_1$;
condition~(ii) forces $\tilde{e}_1 = e_1$.  Inductively, the orthogonal
complement of $\Span(\tilde{e}_1,\ldots,\tilde{e}_{j-1})$ within
$\Span(v_1,\ldots,v_j)$ is one-dimensional, Spanned by~$w_j$.
Normalisation with a positive leading coefficient forces
$\tilde{e}_j = e_j$.
\end{proof}

\begin{example}{Gram--Schmidt in $\R^3$}{ex:gram-schmidt-R3}
\index{Gram--Schmidt process!example}
Apply Gram--Schmidt to $v_1 = (1,1,0)$, $v_2 = (1,0,1)$,
$v_3 = (0,1,1)$ with the standard inner product.

\textbf{Step 1.}\quad $w_1 = v_1 = (1,1,0)$,\quad
$e_1 = w_1/\norm{w_1} = \frac{1}{\sqrt{2}}(1,1,0)$.

\textbf{Step 2.}\quad
$\inner{v_2,e_1} = \frac{1}{\sqrt{2}}(1\cdot 1 + 0\cdot 1 + 1\cdot 0)
  = \frac{1}{\sqrt{2}}$.

$w_2 = v_2 - \inner{v_2,e_1}\,e_1
  = (1,0,1) - \frac{1}{\sqrt{2}}\cdot\frac{1}{\sqrt{2}}(1,1,0)
  = (1,0,1) - (\tfrac{1}{2},\tfrac{1}{2},0)
  = (\tfrac{1}{2},-\tfrac{1}{2},1)$.

$\norm{w_2} = \sqrt{\tfrac{1}{4}+\tfrac{1}{4}+1}
  = \sqrt{\tfrac{3}{2}} = \frac{\sqrt{6}}{2}$,
\quad so $e_2 = \frac{1}{\sqrt{6}}(1,-1,2)$.

\textbf{Step 3.}\quad
$\inner{v_3,e_1} = \frac{1}{\sqrt{2}}(0+1+0) = \frac{1}{\sqrt{2}}$,\quad
$\inner{v_3,e_2} = \frac{1}{\sqrt{6}}(0-1+2) = \frac{1}{\sqrt{6}}$.

$w_3 = v_3 - \frac{1}{\sqrt{2}}\cdot\frac{1}{\sqrt{2}}(1,1,0)
  - \frac{1}{\sqrt{6}}\cdot\frac{1}{\sqrt{6}}(1,-1,2)
  = (0,1,1) - (\tfrac{1}{2},\tfrac{1}{2},0) - (\tfrac{1}{6},-\tfrac{1}{6},\tfrac{1}{3})
  = (-\tfrac{2}{3},\tfrac{2}{3},\tfrac{2}{3})$.

$\norm{w_3} = \tfrac{2}{3}\sqrt{3}$,\quad so
$e_3 = \frac{1}{\sqrt{3}}(-1,1,1)$.

\medskip
The resulting orthonormal basis is:
\[
  e_1 = \frac{1}{\sqrt{2}}\begin{pmatrix}1\\1\\0\end{pmatrix},\quad
  e_2 = \frac{1}{\sqrt{6}}\begin{pmatrix}1\\-1\\2\end{pmatrix},\quad
  e_3 = \frac{1}{\sqrt{3}}\begin{pmatrix}-1\\1\\1\end{pmatrix}.
\]
\end{example}

\begin{center}
\begin{tikzpicture}[>=Stealth, scale=2.2]
  % Gram-Schmidt illustration in 2D (simpler to draw)
  \draw[gray!30, very thin, step=0.5] (-0.3,-0.3) grid (3.3,2.3);
  \draw[->] (-0.3,0) -- (3.3,0) node[right]{$x_1$};
  \draw[->] (0,-0.3) -- (0,2.3) node[above]{$x_2$};

  % v1 and v2
  \draw[->, thick, blue!50, dashed] (0,0) -- (2,1)
    node[right] {\small $v_1$};
  \draw[->, thick, red!50, dashed] (0,0) -- (1,2)
    node[above] {\small $v_2$};

  % e1 (normalised v1)
  \draw[->, very thick, blue!70!black] (0,0) -- ({2/sqrt(5)},{1/sqrt(5)})
    node[below right] {$e_1$};

  % projection of v2 onto e1
  \coordinate (proj) at ({4/(5*sqrt(5))*2},{4/(5*sqrt(5))*1});
  \draw[dotted, gray] (1,2) -- ({4/5},{2/5});

  % w2 = v2 - proj
  \draw[->, thick, orange!80!black] ({4/5},{2/5}) -- (1,2)
    node[midway, left, xshift=-2pt] {\small $w_2$};

  % e2 (normalised w2)
  \pgfmathsetmacro{\wx}{1 - 4/5}
  \pgfmathsetmacro{\wy}{2 - 2/5}
  \pgfmathsetmacro{\wnorm}{sqrt(\wx*\wx + \wy*\wy)}
  \draw[->, very thick, red!70!black] (0,0) --
    ({\wx/\wnorm},{\wy/\wnorm})
    node[above left] {$e_2$};

  % Right angle
  \draw[gray] ({4/5+0.08},{2/5+0.04}) --
    ({4/5+0.04},{2/5+0.12}) -- ({4/5-0.04},{2/5+0.08});
\end{tikzpicture}
\captionof{figure}{Gram--Schmidt in $\R^2$: $v_2$ is orthogonalised
  against $e_1$ to produce $w_2$, then normalised to give~$e_2$.}
\label{fig:gram-schmidt}
\end{center}


% ============================================================
\section{Orthonormal bases}\label{sec:orthonormal-bases}
% ============================================================

\begin{definition}{Orthogonal and orthonormal families}%
{def:orthonormal-family}
\index{orthogonal family}\index{orthonormal family}\index{orthonormal basis}
A family $(e_1,\ldots,e_k)$ in an inner product space is:
\begin{itemize}
  \item \emph{orthogonal} if $\inner{e_i,e_j} = 0$ for all $i\neq j$;
  \item \emph{orthonormal} if it is orthogonal and $\norm{e_i} = 1$ for
    all~$i$.
\end{itemize}
An orthonormal family that is also a basis of~$E$ is called an
\emph{orthonormal basis} (ONB).
\end{definition}

\begin{proposition}{Orthogonal families are linearly independent}%
{prop:orthogonal-indep}
Every orthogonal family of nonzero vectors is linearly independent.
\end{proposition}

\begin{proof}
Let $(e_1,\ldots,e_k)$ be orthogonal with each $e_i\neq\vzero$.
Suppose $\sum_{i=1}^k\alpha_i\,e_i = \vzero$.  Taking the inner
product with~$e_j$:
\[
  0 = \inner*{\sum_i\alpha_i\,e_i,\, e_j}
    = \sum_i\alpha_i\inner{e_i,e_j}
    = \alpha_j\inner{e_j,e_j}
    = \alpha_j\norm{e_j}^2.
\]
Since $\norm{e_j}>0$, we get $\alpha_j = 0$ for all~$j$.
\end{proof}

\begin{corollary}{Existence of orthonormal bases}{cor:onb-existence}
\index{orthonormal basis!existence}
Every finite-dimensional inner product space admits an orthonormal basis.
\end{corollary}

\begin{proof}
Apply the Gram--Schmidt process (
ef{thm:thm:gram-schmidt}) to any basis
of~$E$.
\end{proof}

\begin{theorem}{Parseval's identity}{thm:parseval}
\index{Parseval's identity}
Let $(e_1,\ldots,e_n)$ be an orthonormal basis of an inner product
space~$E$.  Then for all $u,v\in E$:
\[
  \inner{u,v} = \sum_{i=1}^n \inner{u,e_i}\,\overline{\inner{v,e_i}}.
\]
In particular, taking $v = u$:
\[
  \norm{u}^2 = \sum_{i=1}^n \abs{\inner{u,e_i}}^2.
\]
\end{theorem}

\begin{proof}
Write $u = \sum_{i=1}^n\alpha_i\, e_i$ and $v = \sum_{j=1}^n\beta_j\, e_j$
where $\alpha_i = \inner{u,e_i}$ and $\beta_j = \inner{v,e_j}$.  Then
\[
  \inner{u,v} = \inner*{\sum_i\alpha_i\,e_i,\,\sum_j\beta_j\,e_j}
    = \sum_{i,j}\alpha_i\,\bar{\beta}_j\,\inner{e_i,e_j}
    = \sum_i\alpha_i\,\bar{\beta}_i
    = \sum_i\inner{u,e_i}\,\overline{\inner{v,e_i}}. \qedhere
\]
\end{proof}

\begin{theorem}{Bessel's inequality}{thm:bessel}
\index{Bessel's inequality}
Let $(e_1,\ldots,e_k)$ be an orthonormal family (not necessarily a
basis) in an inner product space~$E$.  Then for all $v\in E$:
\[
  \sum_{i=1}^k \abs{\inner{v,e_i}}^2 \leq \norm{v}^2,
\]
with equality if and only if $v\in\Span(e_1,\ldots,e_k)$.
\end{theorem}

\begin{proof}
Let $p = \sum_{i=1}^k\inner{v,e_i}\,e_i$ be the orthogonal projection
of~$v$ onto $F = \Span(e_1,\ldots,e_k)$.  Then $v = p + (v-p)$ with
$p\perp (v-p)$, so by the Pythagorean theorem:
\[
  \norm{v}^2 = \norm{p}^2 + \norm{v-p}^2 \geq \norm{p}^2
    = \sum_{i=1}^k\abs{\inner{v,e_i}}^2.
\]
Equality holds iff $\norm{v-p} = 0$, i.e.\ $v = p\in F$.
\end{proof}

\begin{definition}{Gram matrix}{def:gram-matrix}
\index{Gram matrix}
Let $(v_1,\ldots,v_k)$ be a family of vectors in an inner product
space.  The \emph{Gram matrix}\index{Gram matrix} is the $k\times k$
matrix $G$ defined by
\[
  G_{ij} \deff \inner{v_i,v_j}.
\]
The family is orthonormal if and only if $G = I_k$.  The family is
linearly independent if and only if $\det(G)\neq 0$.
\end{definition}


% ============================================================
\section{Orthogonal and unitary matrices}\label{sec:orthogonal-unitary}
% ============================================================

\begin{definition}{Orthogonal matrix}{def:orthogonal-matrix}
\index{orthogonal matrix}\index{O(n)@$\Orth(n)$}
A matrix $Q\in\Mat_n(\R)$ is \emph{orthogonal} if
\[
  \transpose{Q}\, Q = I_n,
\]
or equivalently $\inv{Q} = \transpose{Q}$.  The set of all $n\times n$
orthogonal matrices is denoted $\Orth(n)$ and forms a group under
matrix multiplication, called the \emph{orthogonal group}.
\end{definition}

\begin{definition}{Unitary matrix}{def:unitary-matrix}
\index{unitary matrix}\index{U(n)@$U(n)$}
A matrix $U\in\Mat_n(\C)$ is \emph{unitary} if
\[
  \hermit{U}\, U = I_n,
\]
or equivalently $\inv{U} = \hermit{U}$.  The set of all $n\times n$
unitary matrices is the \emph{unitary group}~$U(n)$.
\end{definition}

\begin{proposition}{Characterizations of orthogonal/unitary matrices}%
{prop:orthogonal-characterizations}
\index{orthogonal matrix!characterizations}
Let $Q\in\Mat_n(\R)$.  The following are equivalent:
\begin{enumerate}[label=(\roman*)]
  \item $Q$ is orthogonal.
  \item The columns of~$Q$ form an orthonormal basis of~$\R^n$.
  \item The rows of~$Q$ form an orthonormal basis of~$\R^n$.
  \item $Q$ preserves the inner product:
    $\inner{Qx,Qy} = \inner{x,y}$ for all $x,y$.
  \item $Q$ preserves norms: $\norm{Qx} = \norm{x}$ for all~$x$
    (i.e.\ $Q$ is an \emph{isometry}\index{isometry}).
\end{enumerate}
The analogous statements hold for unitary matrices over~$\C$.
\end{proposition}

\begin{proof}
\textbf{(i)$\Leftrightarrow$(ii).}\quad
$\transpose{Q}Q = I$ means $\inner{q_i,q_j} = \delta_{ij}$, where
$q_1,\ldots,q_n$ are the columns of~$Q$.

\textbf{(i)$\Leftrightarrow$(iii).}\quad
$Q\transpose{Q} = I$ (which follows from $\transpose{Q}Q = I$ since
both are square) means the rows are orthonormal.

\textbf{(i)$\Rightarrow$(iv).}\quad
$\inner{Qx,Qy} = \transpose{(Qx)}Qy = \transpose{x}\transpose{Q}Qy
= \transpose{x}y = \inner{x,y}$.

\textbf{(iv)$\Rightarrow$(v).}\quad Take $y = x$.

\textbf{(v)$\Rightarrow$(i).}\quad
$\norm{Qx}^2 = \norm{x}^2$ for all~$x$ means
$\transpose{x}\transpose{Q}Qx = \transpose{x}x$, so
$\transpose{x}(\transpose{Q}Q - I)x = 0$ for all~$x$.  Since
$\transpose{Q}Q - I$ is symmetric, this implies
$\transpose{Q}Q = I$.
\end{proof}

\begin{proposition}{Determinant of orthogonal matrices}%
{prop:det-orthogonal}
If $Q\in\Orth(n)$, then $\det(Q) = \pm 1$.
\end{proposition}

\begin{proof}
$1 = \det(I) = \det(\transpose{Q}Q) = \det(\transpose{Q})\det(Q) =
\det(Q)^2$, so $\det(Q) = \pm 1$.
\end{proof}

\begin{definition}{Special orthogonal group}{def:SO}
\index{SO(n)@$\SO(n)$}\index{special orthogonal group}
The \emph{special orthogonal group} is
$\SO(n) \deff \setcond{Q\in\Orth(n)}{\det(Q) = 1}$.
Its elements are the orthogonal matrices that preserve orientation
(rotations in the geometric sense).
\end{definition}

\begin{example}{Orthogonal matrices in dimension~$2$}{ex:O2}
\index{orthogonal matrix!in dimension 2}
Every matrix in $\SO(2)$ is a rotation:
\[
  R_\theta = \begin{pmatrix} \cos\theta & -\sin\theta \\
    \sin\theta & \cos\theta \end{pmatrix},
  \qquad \theta\in[0,2\pi).
\]
The matrices in $\Orth(2)\setminus\SO(2)$ are reflections:
\[
  S_\theta = \begin{pmatrix} \cos\theta & \sin\theta \\
    \sin\theta & -\cos\theta \end{pmatrix}.
\]
\end{example}

\begin{example}{Permutation matrices}{ex:permutation-orthogonal}
\index{permutation matrix!orthogonality}
Every permutation matrix (obtained by permuting the columns of~$I_n$)
is orthogonal.  Its determinant is $+1$ for even permutations and $-1$
for odd permutations.
\end{example}


% ============================================================
\section{The adjoint operator}\label{sec:adjoint}
% ============================================================

\begin{definition}{Adjoint of a linear map}{def:adjoint}
\index{adjoint operator}\index{adjoint!of a linear map}
Let $(E,\inner{\cdot,\cdot})$ be a finite-dimensional inner product
space and let $f\in\End(E)$.  The \emph{adjoint} of~$f$ is the unique
linear map $f^*\in\End(E)$ satisfying
\[
  \inner{f(u),v} = \inner{u,f^*(v)}
  \qquad\text{for all } u,v\in E.
\]
\end{definition}

\begin{theorem}{Existence and uniqueness of the adjoint}%
{thm:adjoint-existence}
\index{adjoint operator!existence}
For every $f\in\End(E)$ on a finite-dimensional inner product space, the
adjoint $f^*$ exists and is unique.  If $(e_1,\ldots,e_n)$ is an
orthonormal basis and $A = \Mat_{\cB}(f)$ is the matrix of~$f$ in this
basis, then
\[
  \Mat_{\cB}(f^*) = \hermit{A} \quad
  (\text{i.e.\ } \transpose{\bar{A}} \text{ over } \C, \text{ or }
  \transpose{A} \text{ over } \R).
\]
\end{theorem}

\begin{proof}
\textbf{Uniqueness.}\quad Suppose $g,h\in\End(E)$ both satisfy
$\inner{f(u),v} = \inner{u,g(v)} = \inner{u,h(v)}$ for all $u,v$.
Then $\inner{u,g(v)-h(v)} = 0$ for all~$u$, which by non-degeneracy
gives $g(v) = h(v)$ for all~$v$.

\textbf{Existence.}\quad Let $(e_1,\ldots,e_n)$ be an orthonormal basis
and $A = (a_{ij})$ the matrix of~$f$ in this basis, so
$f(e_j) = \sum_i a_{ij}\, e_i$.  Define $g\in\End(E)$ by the matrix
$B = \hermit{A}$, i.e.\ $b_{ij} = \bar{a}_{ji}$.  Then for any two
basis vectors:
\[
  \inner{f(e_j),e_k} = a_{kj},
  \qquad
  \inner{e_j,g(e_k)} = \overline{b_{jk}} = \overline{\bar{a}_{kj}}
    = a_{kj}.
\]
By linearity, $\inner{f(u),v} = \inner{u,g(v)}$ for all $u,v$, so
$f^* = g$ and $\Mat_{\cB}(f^*) = \hermit{A}$.
\end{proof}

\begin{proposition}{Properties of the adjoint}{prop:adjoint-properties}
\index{adjoint operator!properties}
Let $f,g\in\End(E)$ and $\alpha\in\K$.  Then:
\begin{enumerate}[label=(\roman*)]
  \item $(f+g)^* = f^* + g^*$.
  \item $(\alpha f)^* = \bar{\alpha}\, f^*$.
  \item $(f\circ g)^* = g^*\circ f^*$.
  \item $(f^*)^* = f$.
  \item $\Ker(f^*) = (\im f)^{\perp}$.
  \item $\im(f^*) = (\Ker f)^{\perp}$.
\end{enumerate}
\end{proposition}

\begin{proof}
Properties (i)--(iv) follow from the characterizing identity
$\inner{f(u),v} = \inner{u,f^*(v)}$ and uniqueness of the adjoint.

\textbf{(v).}\quad $v\in\Ker(f^*)$ iff $f^*(v) = \vzero$ iff
$\inner{u,f^*(v)} = 0$ for all~$u$ iff $\inner{f(u),v} = 0$ for all~$u$
iff $v\in(\im f)^{\perp}$.

\textbf{(vi).}\quad Apply (v) to $f^*$ to get
$\Ker(f) = (\im f^*)^{\perp}$, then take orthogonal complements:
$\im(f^*) = (\im f^*)^{\perp\perp} = (\Ker f)^{\perp}$.
\end{proof}

\begin{definition}{Self-adjoint, normal, and skew-adjoint operators}%
{def:special-operators}
\index{self-adjoint operator}\index{normal operator}\index{skew-adjoint operator}
An operator $f\in\End(E)$ is called:
\begin{itemize}
  \item \emph{self-adjoint} (or \emph{Hermitian}\index{Hermitian operator})
    if $f^* = f$;
  \item \emph{skew-adjoint} (or \emph{anti-Hermitian}) if $f^* = -f$;
  \item \emph{normal} if $f^*\circ f = f\circ f^*$.
\end{itemize}
In matrix terms (in an ONB): self-adjoint means $A = \hermit{A}$
(symmetric if $\K = \R$), and normal means $\hermit{A}A = A\hermit{A}$.
\end{definition}

\begin{remark}{Orthogonal/unitary operators via the adjoint}%
{rem:orthogonal-adjoint}
An operator $f\in\End(E)$ satisfies $f^*\circ f = \Id$ (equivalently,
$f$ preserves the inner product) if and only if its matrix in an
orthonormal basis is orthogonal (over~$\R$) or unitary (over~$\C$).
In particular, orthogonal and unitary operators are normal.
\end{remark}


% ============================================================
\section{Applications}\label{sec:inner-product-applications}
% ============================================================

% --- Least squares ---

\subsection{Least squares approximation}\label{subsec:least-squares}
\index{least squares}

A fundamental application of orthogonal projection is the \emph{least
squares} method.  Given a system $Ax = b$ that may have no exact
solution (because $b\notin\im(A)$), we seek the vector $\hat{x}$ that
minimises $\norm{Ax - b}$.

\begin{theorem}{Normal equations}{thm:normal-equations}
\index{normal equations}
Let $A\in\Mat_{m\times n}(\R)$ with $m\geq n$ and $\rank(A) = n$.
The least squares solution of $Ax = b$ is the unique vector $\hat{x}$
satisfying
\[
  \transpose{A}\, A\, \hat{x} = \transpose{A}\, b.
\]
The minimum residual is $\norm{b - A\hat{x}} = \norm{b - \proj_{\im A}(b)}$.
\end{theorem}

\begin{proof}
The vector $A\hat{x}$ must be the orthogonal projection of~$b$
onto~$\im(A)$, i.e.\ $b - A\hat{x}\perp\im(A)$.  This means
$\inner{A_j, b - A\hat{x}} = 0$ for each column~$A_j$ of~$A$, which
is equivalent to $\transpose{A}(b - A\hat{x}) = 0$, i.e.\
$\transpose{A}A\hat{x} = \transpose{A}b$.  Since $\rank(A) = n$,
the matrix $\transpose{A}A$ is invertible (it is the Gram matrix of
the columns of~$A$), giving the unique solution
$\hat{x} = \inv{(\transpose{A}A)}\,\transpose{A}\, b$.
\end{proof}

% --- QR factorization preview ---

\subsection{QR factorization}\label{subsec:QR}
\index{QR factorization}

The Gram--Schmidt process has an elegant matrix interpretation.  If
$A\in\Mat_{m\times n}(\R)$ has linearly independent columns
$a_1,\ldots,a_n$, then Gram--Schmidt produces orthonormal vectors
$q_1,\ldots,q_n$ such that
\[
  \Span(a_1,\ldots,a_j) = \Span(q_1,\ldots,q_j)
  \qquad\text{for each } j.
\]
In matrix form, this means $A = QR$, where
\begin{itemize}
  \item $Q\in\Mat_{m\times n}(\R)$ has orthonormal columns
    $q_1,\ldots,q_n$;
  \item $R\in\Mat_n(\R)$ is upper triangular with positive diagonal
    entries, given by $r_{ij} = \inner{a_j,q_i}$ for $i\leq j$.
\end{itemize}
This is the \emph{QR factorization} (or QR decomposition) of~$A$.  It
is numerically more stable than directly solving the normal equations and
is a cornerstone of numerical linear algebra.

% --- Fourier series preview ---

\subsection{Fourier series}\label{subsec:fourier-preview}
\index{Fourier series}

Consider the inner product space $\mathcal{C}([-\pi,\pi],\R)$ with the
$L^2$ inner product $\inner{f,g} = \int_{-\pi}^{\pi}f(t)\,g(t)\,\mathrm{d}t$.
The family
\[
  \frac{1}{\sqrt{2\pi}},\;
  \frac{\cos t}{\sqrt{\pi}},\;
  \frac{\sin t}{\sqrt{\pi}},\;
  \frac{\cos 2t}{\sqrt{\pi}},\;
  \frac{\sin 2t}{\sqrt{\pi}},\;\ldots
\]
is orthonormal.  The orthogonal projection of a function~$f$ onto the
Span of the first $2N+1$ elements is the \emph{truncated Fourier
series}:
\[
  S_N(f)(t) = \frac{a_0}{2} + \sum_{k=1}^{N}\bigl(a_k\cos kt
    + b_k\sin kt\bigr),
\]
where $a_k = \frac{1}{\pi}\int_{-\pi}^{\pi}f(t)\cos kt\,\mathrm{d}t$
and $b_k = \frac{1}{\pi}\int_{-\pi}^{\pi}f(t)\sin kt\,\mathrm{d}t$
are the \emph{Fourier coefficients}\index{Fourier coefficients}.
By 
ef{thm:thm:best-approximation}, this is the best approximation of~$f$
by trigonometric polynomials of degree~$\leq N$ in the $L^2$ norm.
Bessel's inequality (
ef{thm:thm:bessel}) gives
\[
  \frac{a_0^2}{2} + \sum_{k=1}^{N}(a_k^2 + b_k^2)
    \leq \frac{1}{\pi}\int_{-\pi}^{\pi}f(t)^2\,\mathrm{d}t,
\]
and Parseval's identity (
ef{thm:thm:parseval}) asserts that equality
holds in the limit $N\to\infty$ (for $f\in L^2$).


% ============================================================
\section{Exercises}\label{sec:inner-product-exercises}
% ============================================================

\begin{exercise}{Verifying inner product axioms \basic}%
{exo:verify-inner-product}
\index{inner product!verification}
Determine which of the following are inner products on~$\R^2$:
\begin{enumerate}[label=(\alph*)]
  \item $\inner{x,y} = x_1 y_1 - x_1 y_2 - x_2 y_1 + 4 x_2 y_2$.
  \item $\inner{x,y} = x_1 y_1 + x_2 y_2 - x_1 y_2 - x_2 y_1$.
  \item $\inner{x,y} = 2x_1 y_1 + x_1 y_2 + x_2 y_1 + 2 x_2 y_2$.
\end{enumerate}
For each, write the matrix~$G$ and check symmetry and positive
definiteness.
\end{exercise}

\begin{exercise}{Cauchy--Schwarz for sums \basic}{exo:cs-sums}
\index{Cauchy--Schwarz inequality!for sums}
Using the standard inner product on~$\R^n$, deduce from the
Cauchy--Schwarz inequality that for all $a_1,\ldots,a_n,b_1,\ldots,b_n
\in\R$:
\[
  \Bigl(\sum_{i=1}^n a_i b_i\Bigr)^2
    \leq \Bigl(\sum_{i=1}^n a_i^2\Bigr)
         \Bigl(\sum_{i=1}^n b_i^2\Bigr).
\]
When does equality hold?
\end{exercise}

\begin{exercise}{Cauchy--Schwarz for integrals \basic}{exo:cs-integrals}
\index{Cauchy--Schwarz inequality!for integrals}
Deduce from the Cauchy--Schwarz inequality (applied to the $L^2$ inner
product on $\mathcal{C}([a,b])$) that for continuous $f,g$:
\[
  \Bigl(\int_a^b f(t)\,g(t)\,\mathrm{d}t\Bigr)^2
    \leq \int_a^b f(t)^2\,\mathrm{d}t\;\cdot\;
         \int_a^b g(t)^2\,\mathrm{d}t.
\]
\end{exercise}

\begin{exercise}{Orthogonal complement in $\R^3$ \basic}%
{exo:orth-complement-R3}
Let $F = \Span\bigl((1,1,0),\,(0,1,1)\bigr)\subset\R^3$ with the
standard inner product.  Find $F^{\perp}$ and verify that
$\R^3 = F\oplus F^{\perp}$.
\end{exercise}

\begin{exercise}{Gram--Schmidt in $\R^3$ \basic}{exo:gs-R3}
\index{Gram--Schmidt process!exercise}
Apply the Gram--Schmidt process to the basis
$v_1 = (1,0,1)$, $v_2 = (1,1,0)$, $v_3 = (0,1,1)$ of~$\R^3$.
\end{exercise}

\begin{exercise}{Orthogonal projection computation \basic}%
{exo:proj-computation}
In $\R^3$ with the standard inner product, compute the orthogonal
projection of $v = (1,2,3)$ onto the plane
$F = \setcond{(x_1,x_2,x_3)\in\R^3}{x_1 + x_2 + x_3 = 0}$.
\end{exercise}

\begin{exercise}{Orthogonal matrix properties \intermediate}%
{exo:orthogonal-properties}
\index{orthogonal matrix!exercise}
Let $Q\in\Orth(n)$.
\begin{enumerate}[label=(\alph*)]
  \item Show that $\det(Q) = \pm 1$.
  \item Show that if $\lambda$ is an eigenvalue of~$Q$ (possibly
    complex), then $\abs{\lambda} = 1$.
  \item Show that if $n$ is odd and $Q\in\SO(n)$, then $1$ is an
    eigenvalue of~$Q$ (hence $Q$ has a fixed axis).
\end{enumerate}
\end{exercise}

\begin{exercise}{Adjoint of a matrix \intermediate}{exo:adjoint-matrix}
\index{adjoint operator!exercise}
Let $A = \begin{pmatrix} 1 & 2 \\ 3 & 4 \end{pmatrix}\in\Mat_2(\R)$,
viewed as an endomorphism of $(\R^2,\inner{\cdot,\cdot})$ with the
standard inner product.
\begin{enumerate}[label=(\alph*)]
  \item Find $A^* = \transpose{A}$.
  \item Verify the identity $\inner{Ax,y} = \inner{x,A^* y}$ for
    $x = (1,0)$ and $y = (0,1)$.
  \item Is $A$ self-adjoint?  Is $A$ normal?
\end{enumerate}
\end{exercise}

\begin{exercise}{Least squares line fitting \intermediate}%
{exo:least-squares-line}
\index{least squares!line fitting}
Find the line $y = a + bx$ that best fits (in the least squares sense)
the data points $(0,1)$, $(1,0)$, $(2,3)$, $(3,2)$.

\noindent\textit{Hint:} Set up the system $Ax = b$ with
$A = \begin{psmallmatrix} 1 & 0 \\ 1 & 1 \\ 1 & 2 \\ 1 & 3
\end{psmallmatrix}$, $b = (1,0,3,2)^{\mathsf{T}}$, and solve the
normal equations.
\end{exercise}

\begin{exercise}{QR factorization \intermediate}{exo:qr-factorization}
\index{QR factorization!exercise}
Compute the QR factorization of the matrix
\[
  A = \begin{pmatrix} 1 & 1 \\ 1 & 0 \\ 0 & 1 \end{pmatrix}
\]
by applying the Gram--Schmidt process to the columns of~$A$.
Verify that $A = QR$.
\end{exercise}

\begin{exercise}{Isometries of $\R^2$ \intermediate}{exo:isometries-R2}
\index{isometry!classification in $\R^2$}
Show that every $Q\in\Orth(2)$ is either a rotation
$R_\theta\in\SO(2)$ or a reflection.  Describe the geometric action in
each case.
\end{exercise}

\begin{exercise}{Self-adjoint operators have real eigenvalues \intermediate}%
{exo:self-adjoint-eigenvalues}
\index{self-adjoint operator!eigenvalues}
Let $(E,\inner{\cdot,\cdot})$ be a complex inner product space and
$f\in\End(E)$ self-adjoint ($f^* = f$).
\begin{enumerate}[label=(\alph*)]
  \item Show that all eigenvalues of~$f$ are real.
  \item Show that eigenvectors corresponding to distinct eigenvalues
    are orthogonal.
\end{enumerate}
\end{exercise}

\begin{exercise}{Polarization identity \challenging}{exo:polarization}
\index{polarization identity!exercise}
Let $E$ be a real inner product space.  Prove that the inner product is
completely determined by the norm via the \emph{polarization identity}:
\[
  \inner{u,v} = \frac{1}{4}\bigl(\norm{u+v}^2 - \norm{u-v}^2\bigr).
\]
In the complex case, show the analogous identity:
\[
  \inner{u,v} = \frac{1}{4}\sum_{k=0}^{3} i^k\,\norm{u + i^k v}^2.
\]
\end{exercise}

\begin{exercise}{Orthogonal projection onto a line \challenging}%
{exo:proj-line}
Let $a\in\R^n$ be nonzero.  Show that the matrix of the orthogonal
projection onto $\Span(a)$ is
\[
  P = \frac{a\,\transpose{a}}{\transpose{a}\,a}.
\]
Verify that $P^2 = P$, $\transpose{P} = P$, and $\rank(P) = 1$.
\end{exercise}

\begin{exercise}{Unitary diagonalization of Hermitian matrices \challenging}%
{exo:unitary-diag-hermitian}
\index{Hermitian matrix!diagonalization}
Let $A = \begin{pmatrix} 2 & 1-i \\ 1+i & 3 \end{pmatrix}\in\Mat_2(\C)$.
\begin{enumerate}[label=(\alph*)]
  \item Verify that $A$ is Hermitian.
  \item Find the eigenvalues of~$A$ and show they are real.
  \item Find an orthonormal eigenbasis and write $A = U D\hermit{U}$
    for a unitary matrix~$U$ and diagonal matrix~$D$.
\end{enumerate}
\end{exercise}

\begin{exercise}{Fourier coefficients and best approximation \challenging}%
{exo:fourier-best-approx}
\index{Fourier series!exercise}
Let $f(t) = t$ on $[-\pi,\pi]$.
\begin{enumerate}[label=(\alph*)]
  \item Compute the Fourier coefficients $a_k$ and $b_k$ for
    $k = 0,1,2,3$.
  \item Write the best trigonometric polynomial approximation of
    degree~$\leq 3$ in the $L^2$ norm.
  \item Verify Bessel's inequality for $N = 3$.
\end{enumerate}
\end{exercise}


% ============================================================
\section{Chapter summary}\label{sec:inner-product-summary}
% ============================================================

\begin{itemize}
  \item A \textbf{bilinear form} $\varphi\colon E\times E\to\K$ is
    linear in each argument.  It is \textbf{symmetric} if
    $\varphi(u,v) = \varphi(v,u)$ and \textbf{non-degenerate} if its
    matrix is invertible.  The associated \textbf{quadratic form} is
    $q(v) = \varphi(v,v)$.

  \item \textbf{Sylvester's law of inertia}: every real quadratic form
    can be reduced to diagonal form $\diag(+1,\ldots,+1,-1,\ldots,-1,
    0,\ldots,0)$; the \textbf{signature} $(p,r-p)$ is an invariant.

  \item An \textbf{inner product} is a positive definite symmetric
    bilinear form (real case) or a positive definite Hermitian
    sesquilinear form (complex case).

  \item The \textbf{norm} $\norm{v} = \sqrt{\inner{v,v}}$ satisfies
    the \textbf{Cauchy--Schwarz inequality}
    $\abs{\inner{u,v}}\leq\norm{u}\norm{v}$ and the
    \textbf{triangle inequality}
    $\norm{u+v}\leq\norm{u}+\norm{v}$.

  \item Two vectors are \textbf{orthogonal} if $\inner{u,v} = 0$.  The
    \textbf{orthogonal complement} $F^{\perp}$ satisfies
    $E = F\oplus F^{\perp}$.

  \item The \textbf{orthogonal projection} $\proj_F(v)$ is the unique
    closest point to~$v$ in a subspace~$F$ (\textbf{best approximation
    property}).

  \item The \textbf{Gram--Schmidt process} turns any linearly
    independent family into an orthonormal one while preserving the
    successive Spans.  Every inner product space has an
    \textbf{orthonormal basis}.

  \item \textbf{Parseval's identity}: in an ONB,
    $\inner{u,v} = \sum_i\inner{u,e_i}\overline{\inner{v,e_i}}$.
    \textbf{Bessel's inequality}: partial sums of Fourier coefficients
    are bounded by $\norm{v}^2$.

  \item \textbf{Orthogonal matrices} ($\transpose{Q}Q = I$) preserve
    inner products and norms; they form the group $\Orth(n)$, with the
    subgroup $\SO(n)$ (rotations, $\det = 1$).  Over~$\C$, the
    analogues are \textbf{unitary matrices}.

  \item The \textbf{adjoint} $f^*$ of a linear map satisfies
    $\inner{f(u),v} = \inner{u,f^*(v)}$; in an ONB, its matrix is
    $\hermit{A}$.  Key identities:
    $\Ker(f^*) = (\im f)^{\perp}$ and
    $\im(f^*) = (\Ker f)^{\perp}$.

  \item \textbf{Applications}: \emph{least squares} (solve overdetermined
    systems via $\transpose{A}A\hat{x} = \transpose{A}b$), \emph{QR
    factorization} (matrix form of Gram--Schmidt), and \emph{Fourier
    series} (orthogonal projection onto trigonometric polynomials).
\end{itemize}
